\documentclass[a4paper]{article}
\usepackage[pdftex]{graphicx}
\usepackage[utf8]{inputenc}
\usepackage{enumerate}
\usepackage{amssymb}
\usepackage{icomma}
\usepackage{siunitx}
\sisetup{locale=DE}
\usepackage{href-ul}
\hypersetup{
	colorlinks=true,
	linkcolor=blue,
	urlcolor=blue}
\usepackage{geometry}
\geometry{a4paper, top=15mm, left=15mm, right=15mm, bottom=15mm,
	headsep=10mm, footskip=12mm}

\begin{document}
	
	{\bf Horizontal throw}
	\begin{enumerate}[1.]
		\begin{minipage}{0.6\textwidth}
			{\bf \item Task }
			\begin{enumerate}[a)]
				\item In the arrangement shown below there is a ball in position A. A is twice as high above the ground as B.
				Now the ball is dropped.
				It bounces elastically horizontally at point B and finally lands at point C
				on the floor. Calculate the throw distance $x$ for $y=\SI{1,0}{\meter}$.
				\item Carry out the same calculation again in general and represent the width $x(y)$ as a function of $y$.
				\item Does the throw distance change if this arrangement is operated on the moon $g_{Moon}=\SI{1,6}{\meter \over {\second^2}}$?
				Reasons!
			\end{enumerate}
		\end{minipage}
		\hspace{0.5 cm}
		\begin{minipage}{0.4\textwidth}
			\includegraphics[width=5 cm]{wamm074}
		\end{minipage}
		
		\begin{minipage}{0.6\textwidth}
			{\bf \item A firefighting plane has filled up with a load of water} and is now flying at a speed of $v=\SI{180}{\kilo\meter \over \hour}$
			at a height of $h=\SI{200}{\meter}$ above the forest. The source of the fire is still small and must be hit precisely. At what horizontal distance?
			Does the plane have to open the flaps so that the water hits the source of the fire?
			{\bf \item Sports shooter Lunerald is sure} that he aimed at the center of the target. Nevertheless, the shot missed. His rifle accelerated the bullet to $v=\SI{310}{\meter\over\second}$. How far was the shooting distance if the second ring of the target has a radius of $\SI{2.0}{\centi\meter}$?
		\end{minipage}
		\hfill
		\begin{minipage}{0.4\textwidth}
			\includegraphics[width=5 cm]{suss074}
		\end{minipage}
		
		\begin{minipage}{0.6\textwidth}
			{\bf \item A golf ball has a diameter of} $\SI{43}{\milli\meter}$, the hole diameter is $\SI{13.4}{\centi\meter}$. A golf ball will be properly "sunk" if it is not too fast. Otherwise it will bounce up on the edge of the hole.
			\begin{enumerate}[a)]
				\item What horizontal and vertical distance and does it fall from the edge of the hole if its "equator" should touch the opposite edge?
				\item Calculate the maximum speed that the ball can have to land in the hole.
			\end{enumerate}
		\end{minipage}
		\hspace{0.5 cm}
		\begin{minipage}{0.4\textwidth}
			\includegraphics[width=5 cm]{golb074}
		\end{minipage}
		
		\begin{minipage}{0.6\textwidth}
			{\bf \item Spiff, the cliff diver takes a running start...} What is the minimum horizontal speed that his jump must have so that he falls safely over the inclined cliff into the water? Here the height is $h=\SI{14}{\meter}$ and to be on the safe side it should be $\SI{2.0}{\meter}$ next to the edge of the rock.
			\vspace{1cm}
			
			\fbox{
				\begin{minipage}{0.4\textwidth}
					For the solution, please \href{https://www.okuyakl.de/phys/p9wawL074/le074.pdf}{click here} or scan the QR code.\\
					Additional worksheets are available at
					
					\href{http://www.okuyakl.com}{www.okuyakl.com}
				\end{minipage}
				\hfill
				\begin{minipage}{0.55\textwidth}
					\includegraphics[width=3 cm]{qre074}
					\includegraphics[width=2 cm]{../../viecher/afanticon1}
					
			\end{minipage}}
			
		\end{minipage}
		\hfill
		\begin{minipage}{0.4\textwidth}
			\includegraphics[width=5 cm]{cliff074}
		\end{minipage}
		
	\end{enumerate}
	
	
\end{document}